% Autor: Elias Welt
\section{Projektsetup und Konfiguration}
Das Projekt wurde mit dem Befehl \texttt{npm create vite@latest} initialisiert. Die Konfiguration befindet sich in der Datei \texttt{vite.config.js}.

\begin{lstlisting}[language=JavaScript, caption={Auszug aus der Vite-Konfiguration}]
import { defineConfig } from 'vite'
import react from '@vitejs/plugin-react'

// https://vite.dev/config/
export default defineConfig({
  plugins: [react()],
  server: {
    host: true,
    port: 5173,
    proxy: {
      '/api': {
        target: process.env.VITE_API_TARGET || 'http://localhost:5015',
        changeOrigin: true
      }
    }
  }
})
\end{lstlisting}

Ein kritischer Aspekt dieser Konfiguration ist der \textbf{Proxy}: Um während der Entwicklung Cross-Origin Resource Sharing (CORS) Probleme zu vermeiden, leitet der Vite-Entwicklungsserver alle Anfragen, die mit \texttt{/api} beginnen, an das Backend (Port 5015) weiter. Die Option \texttt{changeOrigin: true} manipuliert dabei den Host-Header der Anfrage so, als käme sie direkt vom Backend-Server selbst.

Die Projektstruktur im Ordner \texttt{src} folgt einer funktionalen Aufteilung:
\begin{itemize}
    \item \texttt{components/}: Beinhaltet UI-Komponenten (z.B. \texttt{Chat.jsx}, \texttt{Quiz.jsx}). Jede Komponente hat idealerweise ihre eigene CSS-Datei (z.B. \texttt{Chat.css}) für Styling
    \item \texttt{services/}: Inkludiert die gesamte API-Kommunikation
    \item \texttt{assets/}: Für statische Dateien
\end{itemize}
